% TOP_PROBLEM: {"id": 1164, "title": "Falconer's Conjecture", "category_id": 6, "category": {"id": 6, "name": "geometry", "display_name": "Geometry", "description": "Euclidean and non-Euclidean geometry, geometric structures.", "slug": "geometry", "order_index": 6, "created_at": "2026-07-31T15:26:25.671Z"}, "status": "open"}
% FIRST_POSTED: null
% ENTRY_KIND: statement_only
% Imported from: batch04.tex; UnsolvedMath ID 1164
\documentclass[11pt]{article}
\usepackage[margin=25mm]{geometry}
\usepackage{fontspec}
\setmainfont{Latin Modern Roman}
\usepackage{amsmath,amssymb,mathtools}
\usepackage{unicode-math}
\setmathfont{Latin Modern Math}
\usepackage{xurl,hyperref}
\hypersetup{colorlinks=true,urlcolor=blue}
\setcounter{secnumdepth}{0}
\setlength{\parindent}{0pt}
\setlength{\parskip}{5pt}
\setlength{\emergencystretch}{3em}
\newcommand{\sref}[2]{\hyperlink{source:#1:#2}{[S#2]}}
\newcommand{\cataloguescope}[1]{\paragraph{Recorded target.}#1}
\begin{document}
% UnsolvedMath ID 1164; source catalogue code GEO-009
\setcounter{section}{1163}
\section{Falconer's Conjecture}\label{1164}
{\small\sffamily UnsolvedMath ID 1164 · Geometry}
\subsection{Definitions and mathematical statement}

\paragraph{Shared notation.}$\mathbb N=\{1,2,\ldots\}$, and $\mathbb Z,\mathbb Q,\mathbb R,\mathbb C$ denote the integers, rationals, reals and complex numbers. Any other conventions are specified locally.


For a compact set \(E\subseteq\mathbb R^d\), its Euclidean
distance set is
\[
 \Delta(E)=\{\|x-y\|:x,y\in E\}\subseteq[0,\infty).
\]
Write \(\mathcal L^1\) for one-dimensional Lebesgue measure.
The Hausdorff dimension \(\dim_H E\) is
\(\inf\{s\ge0:\mathcal H^s(E)=0\}\), where
\[
 \mathcal H^s(E)=
 \lim_{\delta\downarrow0}\inf
 \left\{\sum_i(\operatorname{diam}U_i)^s:
       E\subseteq\bigcup_iU_i,\ \operatorname{diam}U_i<\delta\right\}.
\]
An inessential positive normalization factor in \(\mathcal H^s\)
does not change this dimension.
\paragraph{Statement recorded in the supplied source.}
For every integer \(d\ge2\) and every compact set
\(E\subseteq\mathbb R^d\), does
\[
 \dim_H E>\frac d2
       \quad\Longrightarrow\quad
 \mathcal L^1(\Delta(E))>0
\]
hold?  The same implication in dimension one is also included by
the literal dimension-free wording, but the higher-dimensional
assertion above is the standard Falconer problem. \sref{1164}{2} \sref{1164}{3}
\subsection{Short English statement}
Does every compact set of Hausdorff dimension greater than half the ambient dimension determine a positive-measure set of distances?
\subsection{Sources}
\begin{description}
\item[\hypertarget{source:1164:1}{S1}] ULAM.AI and the UnsolvedMath Contributors, \emph{UnsolvedMath: A Curated Collection of Open Mathematics Problems}, record 1164 (GEO-009). CC BY 4.0. \url{https://huggingface.co/datasets/ulamai/UnsolvedMath}.
\item[\hypertarget{source:1164:2}{S2}] Larry Guth, Alex Iosevich, Yumeng Ou and Hong Wang, \emph{On Falconer's distance set problem in the plane}, Inventiones Mathematicae 219 (2020), 779--830, Section 1; arXiv:1808.09346. \url{https://arxiv.org/abs/1808.09346}.
\item[\hypertarget{source:1164:3}{S3}] Xiumin Du, Alex Iosevich, Yumeng Ou, Hong Wang and Ruixiang Zhang, \emph{An improved result for Falconer's distance set problem in even dimensions}, Mathematische Annalen 380 (2021), 1215--1231, Section 1; arXiv:2006.06833. \url{https://arxiv.org/abs/2006.06833}.
\end{description}
\label{end:1164}
\end{document}
